% s-231-p231-m001; original PDF 231, 232
\Aphorism{s-231-p231-m001}{326}
Kant hielt grosse Stücke auf seine Kategorientafel, Schopenhauer nannte sich einen Oligographen, Nietzsche rühmt sich seiner ergangenen, nicht ersessenen Gedanken. Von Künstlern ist man nichts Anderes als Selbst-Missverständniss gewöhnt; aber dass selbst Philosophen durch den Intellect, der ihre Hauptwaffe ist, nicht davor geschützt sind, in einem Zufälligen oder auch nur Eingebildeten ihr Wesentliches zu sehen — das giebt zu erstaunen. Kants Dodekalog ist etwa zu den Fällen associativen Zwanges zu zählen, die in der Geschichte wissenschaftlicher Darstellung eine so grosse Rolle spielen: der neue Wein will zum alten Schlauch, — der Entdecker einer Wahrheit vermag nicht, diese Wahrheit ihre eigene Sprache reden zu lassen, weil ihm die Formalia älterer Wahrheiten als „ausgefahrene Geleise“ im Gehirn sitzen. Der Erkenntnisskritiker Kant konnte nicht vergessen, dass er eigentlich Professor der Logik war. Schopenhauer mag, mit seinen sechs Bänden, im Verhältniss zu zwölf- und achtzehnbändigen „Universitätsphilosophen“ eine Art Oligograph sein; in sich selbst vertrüge er eine weitere Zusammenziehung auf ein Sechstel seines Volumens — wohlgemerkt, ich rede von seiner Gedankensubstanz, nicht von der Ausführung, vom Forscher, nicht vom Dichter Schopenhauer. Über Nietzsches Production sind wir durch die begonnene Herausgabe seines Nachlasses unterrichtet genug, die Weite der Distanz, die Energie von Schreibtischarbeit zwischen dem ersten Auftauchen des „ergangenen“ Gedankens und seiner letzten druckfertigen Niederschrift verfolgen zu können; ohne das „Sitzfleisch“, das Nietzsche die Sünde wider den heiligen Geist nennt, wäre es auch hier nicht abgegangen. Wenn Feder und Papier aber auf den Zwischenstufen der Gedankenschöpfung thätig sind, so ist es blosse Velleität, sie als Geburtshelfer oder Förderer der ersten Conception zu verwerfen. Dass der Eine gehend, der Andere sitzend oder liegend am besten denkt, drückt nur eine Relation zwischen sensiblen und motorischen Nerven aus, keinen Rangunterschied des Denkens selbst; dass Dieser auf eine Vormittagslandschaft, Jener auf Kerzenlicht, ein Dritter auf einen Bogen weissen Papiers blicken muss, um fruchtbar zu sein, besagt nichts über Kraft und Saft der Früchte. Diese Anreger und Auslöser der Gehirnthätigkeit sind ja nur Gelegenheitsursachen; die eigentliche Ursache hat lange vorher, mit ganz anderen Wonnen und Wärmen gewirkt.
\par

